\section*{第 \ref{ch:b3:complex-mechanisms} 章——配合物的反应机理}

\begin{solution}{exo:b3:complex-mechanisms:1}
惰性：$\ce{[Cr(H2O)6]^{3+}}$（$d^3$）和 $\ce{[Co(NH3)6]^{3+}}$（低自旋 $d^6$），配体场
活化能较大。活性的：$\ce{[Cr(H2O)6]^{2+}}$（高自旋 $d^4$）和 $\ce{[Cu(H2O)6]^{2+}}$（$d^9$），
发生姚–泰勒畸变，轴向键长而弱；$\ce{[Zn(H2O)6]^{2+}}$（$d^{10}$），无配体场
势垒。
\end{solution}

\begin{solution}{exo:b3:complex-mechanisms:2}
当 $k_2[\mathrm Y] \gg k_{-1}[\mathrm X]$ 时，分母为 $k_2[\mathrm Y]$，$v = k_1[\mathrm{ML_5X}]$；当 $k_{-1}[\mathrm X] \gg k_2[\mathrm Y]$ 时，
分母为 $k_{-1}[\mathrm X]$，由此得第二式（
即由离去配体抑制的预平衡）。
\end{solution}

\begin{solution}{exo:b3:complex-mechanisms:3}
D（解离）：$\Delta V^{\ddagger} > 0$，$\Delta^{\ddagger}S^\circ > 0$（配体离去）。A（缔合）：两者均为负（
过渡态中有配体结合）。
\end{solution}

\begin{solution}{exo:b3:complex-mechanisms:4}
$cis$-$\ce{[PtCl2(NH3)2]}$ 由 $\ce{[PtCl4]^{2-}}$ 制得；$trans$-$\ce{[PtCl2(NH3)2]}$ 由 $\ce{[Pt(NH3)4]^{2+}}$ 制得（
\ce{Cl-} 的\omterm{def:b3:complex-mechanisms:trans-effect}{反位效应}大于 \ce{NH3}）。
\end{solution}

\begin{solution}{exo:b3:complex-mechanisms:5}
$k_{\mathrm{obs}} = 2.0 \times \num{3.0e4} \times 0.10/1.2 = \qty{5.0e3}{s^{-1}}$；在 \qty{1.0}{mol/L} 时，$\num{6.0e4}/3.0 = \qty{2.0e4}{s^{-1}}$；极限
值 $k_i = \qty{3.0e4}{s^{-1}}$。
\end{solution}

\begin{solution}{exo:b3:complex-mechanisms:6}
$\Delta V^{\ddagger} = -RT\ln2/\Delta p = -8.314 \times 298 \times 0.693/\num{99.9e6} = \qty{-1.7e-5}{m^3.mol^{-1}} = \qty{-17}{cm^3.mol^{-1}}$：
\omterm{def:b3:complex-mechanisms:mechanisms}{缔合机理}。
\end{solution}

\begin{solution}{exo:b3:complex-mechanisms:7}
膦是强 $\sigma$ 给体、$\pi$ 受体，削弱了与其反位的键（\omterm{def:b3:complex-mechanisms:trans-effect}{反位
影响}，键变长），并使其活化（\omterm{def:b3:complex-mechanisms:trans-effect}{反位效应}）：该氯离子
最先被取代。
\end{solution}

\begin{solution}{exo:b3:complex-mechanisms:8}
桥联配体最终转移到被氧化的金属上；速率
强烈依赖于潜在桥联配体的性质；检测到双核中间体；
且速率不超过在活性伙伴上形成桥
所需取代反应的速率。
\end{solution}

\begin{solution}{exo:b3:complex-mechanisms:9}
$(\lambda + \Delta G^\circ)^2/4\lambda$，其中 $4\lambda = 320$：20.0、11.3、0 和 \qty{5.0}{kJ/mol}（最后一个
位于\omterm{def:b3:complex-mechanisms:inverted}{反转区}）。
\end{solution}

\begin{solution}{exo:b3:complex-mechanisms:10}
$k_{12} = \sqrt{4.0 \times \num{2.0e5} \times \num{1.0e4}} = \qty{8.9e4}{L.mol^{-1}.s^{-1}}$。
\end{solution}

\begin{solution}{exo:b3:complex-mechanisms:11}
高自旋 $d^5$：0；$d^7$：$1.33\,Dq$；$d^8$：$2.0\,Dq$。预期的水交换速率：\ce{Mn^{2+}} $>$ \ce{Co^{2+}} $>$
\ce{Ni^{2+}}，与观察到的顺序一致。
\end{solution}

\begin{solution}{exo:b3:complex-mechanisms:12}
对于强放能反应，本征速率超过了离子
相遇的速率，观测速率常数保持在扩散极限，下降
被掩盖。当供体与受体固定于同一分子中的一定距离时，
不存在扩散步骤；一系列驱动力递增的受体显示
速率先升高、经过极大值后下降。
\end{solution}

\begin{solution}{pb:b3:complex-mechanisms:1}
\textbf{1.} $5d^8$。
\textbf{2.} $5d$ 离子的强场使指向四个
配体的 $d_{x^2-y^2}$ 轨道能级很高：八个电子在四个较低轨道中成对。
\textbf{3.} 该配合物有 16 个价电子，轴向有一个空位：进入的
配体先结合，经五配位过渡态。
\textbf{4.} $k_{\mathrm{obs}} = k_1 + k_2[\mathrm Y]$：溶剂途径（溶剂进攻，随后被 Y 快速取代）
和直接途径。
\textbf{5.} 取代需要数小时（见第三部分）：以分钟计为惰性，
而不是以天计。
\textbf{6.} 得到顺铂，因为第二个 \ce{NH3} 取代的是与氯处于反位的氯，
（\omterm{def:b3:complex-mechanisms:trans-effect}{反位效应} $\ce{Cl-} > \ce{NH3}$）；同时也生成副产物。
\textbf{7.} 碘离子的\omterm{def:b3:complex-mechanisms:trans-effect}{反位效应}大于氯离子，这使得定位
效应更清晰、更快。
\textbf{8.} 在 $\ce{[PtI3(NH3)]-}$ 中，与碘离子处于反位的碘离子被活化，而不是与
\ce{NH3} 处于反位的那个；第二个氨分子进入时与第一个呈顺式。
\textbf{9.} 碘化银沉淀，$cis$-$\ce{[Pt(H2O)2(NH3)2]^{2+}}$ 仍留在溶液中。
\textbf{10.} 氯离子在各自原位取代两个水分子；
氨配体不受影响。
\textbf{11.} 由 $\ce{[Pt(NH3)4]^{2+}}$ 与两个氯离子制得。
\textbf{12.} 其两个活性位置彼此相对，
不能同时结合同一 DNA 链上相邻的两个碱基，
这正是顺铂发挥作用的损伤方式。
\textbf{13.} $\ln2/\num{9.0e-5} = \qty{7.7e3}{s} = \qty{2.1}{h}$。
\textbf{14.} $1 - \eu^{-\num{9.0e-5} \times 7200} = 0.48$。
\textbf{15.} 水作为离去基团远优于氯离子：水合配合物
很快被 DNA 碱基的氮取代。
\textbf{16.} \omterm{def:b3:complex-mechanisms:mechanisms}{缔合机理}，与所有铂(II)取代反应一样：\omterm{def:b3:complex-mechanisms:activation-volume}{活化
体积}为负。
\textbf{17.} 高氯离子浓度使平衡向回到
二氯形式的方向移动，使药物到达细胞之前保持完整。
\textbf{18.} $f = K/(K + [\ce{Cl-}])$。
\textbf{19.} $\num{3.0e-3}/0.103 = 0.029$。
\textbf{20.} $\num{3.0e-3}/0.0070 = 0.43$。
\textbf{21.} 15。
\textbf{22.} 只有在细胞内，相当大的比例才以活性水合
配合物形式存在。
\textbf{23.} 含硫配体，如谷胱甘肽以及蛋白质中的半胱氨酸和蛋氨酸，
它们能与软性的铂强烈结合。
\textbf{24.} 氨与铂结合很强，位于与氯处于反位的位置，是\omterm{def:b3:complex-mechanisms:trans-effect}{反位效应}较小的配体，
因此氨配体的键不会被活化。
\textbf{25.} \textbf{细胞内以活性水合形式存在的药物量
约为血浆中的 15 倍。}
\end{solution}
